%% ma: L11-B17-0068
%% lop: 11
%% chuong: 5
%% bai: 17
%% dang: 11.17.4
%% loai: DS
%% muc: [TH, TH, TH, TH]
%% dap_an: "SSĐS"
%% nguon: "tai-lieu-hoc-tap-mon-toan-lop-11-theo-chuong-trinh-giao-duc-pho-thong-2018.pdf (Lớp Toán Thầy Duy, Phạm Lê Duy, THPT Chu Văn An), Chương 3 Bài 3, Phần 2, Câu 8 (tr.417)"
%% kien_thuc: "Lớp 11 Bài 17 (hàm số liên tục)"
%% trang_thai: chinh-thuc
%% ghi_chu: "Đã duyệt 03/11/2026."
\begin{ex}
Cho hàm số $f(x)=\begin{cases}\dfrac{2x+6}{3x^2-27}&\text{khi }x\ne\pm3\\-\dfrac19&\text{khi }x=\pm3\end{cases}$. Khi đó
\choiceTF
{Hàm số liên tục tại mọi điểm trừ các điểm thuộc khoảng $(-3;3)$.}
{Hàm số liên tục tại mọi điểm trừ điểm $x=-3$.}
{\True Hàm số liên tục tại mọi điểm trừ điểm $x=3$.}
{Hàm số liên tục trên $\mathbb R$.}
\loigiai{\sloppy Với $x\ne\pm3$: $f(x)=\dfrac{2(x+3)}{3(x-3)(x+3)}=\dfrac2{3(x-3)}$ liên tục. Tại $x=-3$: $\lim\limits_{x\to-3}f(x)=\dfrac2{3(-6)}=-\dfrac19=f(-3)$, liên tục. Tại $x=3$: $\lim\limits_{x\to3}\dfrac2{3(x-3)}$ không hữu hạn ($f$ không có giới hạn hữu hạn) nên gián đoạn tại $x=3$.
Vậy hàm số liên tục tại mọi điểm trừ $x=3$: chỉ ý c) đúng.}
\end{ex}
